
\documentclass{article}
\usepackage{colortbl}
\usepackage{makecell}
\usepackage{multirow}
\usepackage{supertabular}

\begin{document}

\newcounter{utterance}

\centering \large Interaction Transcript for game `cladder', experiment `full\_v1.5\_default', episode 6304 with llama{-}3.1{-}8b{-}instruct.
\vspace{24pt}

{ \footnotesize  \setcounter{utterance}{1}
\setlength{\tabcolsep}{0pt}
\begin{supertabular}{c@{$\;$}|p{.15\linewidth}@{}p{.15\linewidth}p{.15\linewidth}p{.15\linewidth}p{.15\linewidth}p{.15\linewidth}}
    \# & \multicolumn{2}{c}{Player} && \multicolumn{2}{c}{Game Master} \\
    \hline

    \theutterance \stepcounter{utterance}  
    & & & \multicolumn{4}{p{0.6\linewidth}}{
        \cellcolor[rgb]{0.9,0.9,0.9}{
            \makecell[{{p{\linewidth}}}]{
                \texttt{\tiny{[P1$\langle$GM]}}
                \texttt{You are an expert in causal inference. The following question is not a typical commonsense query, but rather a meticulously designed question created by a professor specializing in causal inference, intended to assess the students' mastery of the course content.} \\
\texttt{Imagine a self{-}contained, hypothetical world with only the following conditions, and without any unmentioned factors or causal relationships: Zuph has a direct effect on rixq and xevu. Rixq has a direct effect on xevu. The overall probability of rixq is 32\%. The probability of not rixq and xevu is 35\%. The probability of rixq and xevu is 21\%. Is the chance of xevu smaller when observing rixq?} \\
\texttt{Start your answer with 'Yes' or 'No', followed by additional reasoning or evidence to support your explanation.} \\
            }
        }
    }
     \\ \\

    \theutterance \stepcounter{utterance}  
    & \multicolumn{4}{p{0.6\linewidth}}{
        \cellcolor[rgb]{0.561,0.737,0.561}{
            \makecell[{{p{\linewidth}}}]{
                \texttt{\tiny{[P1$\rangle$GM]}}
                \texttt{No.} \\
\\ 
\texttt{To determine if the chance of xevu is smaller when observing rixq, we need to calculate the conditional probability of xevu given rixq (P(xevu|rixq)) and compare it to the unconditional probability of xevu (P(xevu)).} \\
\\ 
\texttt{We can use the following formula to calculate the conditional probability:} \\
\\ 
\texttt{P(xevu|rixq) = P(rixq and xevu) / P(rixq)} \\
\\ 
\texttt{Given values:} \\
\texttt{{-} P(rixq and xevu) = 21\%} \\
\texttt{{-} P(rixq) = 32\%} \\
\\ 
\texttt{P(xevu|rixq) = 21\% / 32\% ≈ 0.65625} \\
\\ 
\texttt{Now, we need to calculate the unconditional probability of xevu (P(xevu)). We can use the fact that the probability of not rixq and xevu is 35\%. Since the probability of not rixq is 1 {-} P(rixq) = 1 {-} 32\% = 68\%, we can calculate P(xevu) as follows:} \\
\\ 
\texttt{P(xevu) = P(not rixq and xevu) + P(rixq and xevu)} \\
\texttt{= 35\% + 21\%} \\
\texttt{= 56\%} \\
\\ 
\texttt{Now, we can compare P(xevu|rixq) and P(xevu). Since P(xevu|rixq) ≈ 0.65625 is greater than P(xevu) = 56\%, the chance of xevu is actually larger when observing rixq, not smaller.} \\
            }
        }
    }
    & & \\ \\

    \theutterance \stepcounter{utterance}  
    & & & \multicolumn{2}{p{0.3\linewidth}}{
        \cellcolor[rgb]{0.9,0.9,0.9}{
            \makecell[{{p{\linewidth}}}]{
                \texttt{\tiny{[GM$|$GM]}}
                \texttt{No} \\
            }
        }
    }
    & & \\ \\

    \theutterance \stepcounter{utterance}  
    & & & \multicolumn{2}{p{0.3\linewidth}}{
        \cellcolor[rgb]{0.9,0.9,0.9}{
            \makecell[{{p{\linewidth}}}]{
                \texttt{\tiny{[GM$|$GM]}}
                \texttt{no} \\
            }
        }
    }
    & & \\ \\

    \theutterance \stepcounter{utterance}  
    & & & \multicolumn{2}{p{0.3\linewidth}}{
        \cellcolor[rgb]{0.9,0.9,0.9}{
            \makecell[{{p{\linewidth}}}]{
                \texttt{\tiny{[GM$|$GM]}}
                \texttt{game\_result = WIN} \\
            }
        }
    }
    & & \\ \\

\end{supertabular}
}

\end{document}
